% --- File: sections/01_introduction.tex ---

Modern institutional portfolios face unprecedented challenges in managing tail risk amidst structural macroeconomic regime shifts. As traditional asset classes increasingly exhibit non-stationary dynamics, severe volatility clustering, and fat-tailed return distributions, legacy risk frameworks---ranging from Historical Simulation to parametric Variance-Covariance approaches---frequently fail to rapidly adapt, leaving institutions vulnerable to catastrophic tail events.

Under the impending regulatory mandates of Basel IV and the Fundamental Review of the Trading Book (FRTB), the accuracy of internal models is paramount. Financial institutions face punitive capital add-ons if their 99\% Value-at-Risk (VaR) models experience frequent breaches during backtesting. To avoid regulatory penalties and forced migration to the Standardized Approach (SA), risk managers often over-calibrate models, leading to artificially inflated capital reserves and severe capital inefficiency. 

In recent years, Deep Learning---specifically Recurrent Neural Networks (RNNs) and Long Short-Term Memory (LSTM) architectures---has been proposed to capture the complex, non-linear dynamics of financial time series. However, deploying pure ``black-box'' machine learning in risk production environments introduces a new set of critical vulnerabilities. Naive LSTMs struggle to isolate the true underlying data generation process in highly noisy environments, frequently resulting in severe overfitting during low-volatility regimes and catastrophic under-prediction during structural market breaks.

To bridge the gap between adaptive machine learning and robust regulatory compliance, this paper introduces an ``Anchored'' Neural Architecture for generalized tail-risk estimation. Drawing inspiration from Physics-Informed Deep Learning, we propose a hybrid framework that constrains the neural network using a rolling parametric statistical prior. By utilizing an asymmetric Quantile Loss function (Pinball Loss) anchored to a rolling parametric VaR estimate, the model explicitly targets the 1st percentile of the return distribution while maintaining the statistical safety bounds required by regulators.

To rigorously evaluate the resilience of this architecture, we utilize high-frequency digital asset data (BTC/USD) not as a target investment, but as an extreme stress-testing proxy. Assets with such severe leptokurtic distributions provide a worst-case validation environment; an architecture capable of surviving these extreme dynamics will seamlessly handle the volatility of traditional Tier-1 equities or FX portfolios.

The primary contribution of this research is twofold. First, we demonstrate mathematically and empirically why unconstrained neural networks fail during out-of-sample stress tests. Second, we present a generalized, potentially production-relevant methodology that materially reduces required regulatory capital---freeing up \capEfficiencyBps{} of capital reserves in our proxy environment---without compromising the strict breach-rate limits mandated by the Kupiec Proportion of Failures test.

The remainder of this paper is structured as follows. Section 2 outlines the regulatory context of Basel IV/FRTB. Section 3 details the methodology, comparing the Historical VaR baseline, the Naive LSTM, and our proposed Anchored architecture. Section 4 describes the experimental setup and the stress-testing framework. Finally, Section 5 presents the empirical results, risk validation metrics, and implications for institutional capital efficiency.